                \begin{tikzpicture}[xscale=1.0, yscale=1.0]
                    \useasboundingbox (-1.05, -1.25) rectangle (14.15, 4.75);

                    \tikzset{axiome/.style={draw=gris_fonce_inria, line width=0.9pt, rounded corners=3pt,
                                            fill=gray!6, inner sep=4pt, align=center, font=\scriptsize}}
                    \tikzset{coupe/.style={inria-rouge, densely dashed, line width=0.8pt}}
                    \tikzset{etiq/.style={font=\scriptsize, text=inria-rouge, inner sep=1.5pt}}

                    % ---------- Panneau gauche : les trois axiomes ----------
                    \node[axiome] (ax1) at (1.60, 3.85)
                        {\textsf{(i)} invariance d'échelle\\{\tiny $d \mapsto \lambda\,d$ : même résultat}};
                    \node[axiome] (ax2) at (0.05, 1.55)
                        {\textsf{(ii)} richesse\\{\tiny toute partition est atteignable}};
                    \node[axiome] (ax3) at (3.20, 1.55)
                        {\textsf{(iii)} cohérence\\{\tiny resserrer un amas ne change rien}};

                    \draw[gris_fonce_inria, line width=0.7pt] (ax1) -- (ax2);
                    \draw[gris_fonce_inria, line width=0.7pt] (ax2) -- (ax3);
                    \draw[gris_fonce_inria, line width=0.7pt] (ax3) -- (ax1);

                    \node[font=\Large\bfseries, text=inria-rouge] at (1.62, 2.62) {$\nexists$};

                    \node[font=\scriptsize, align=center, text=gris_fonce_inria] at (1.62, 0.50)
                        {aucune \emph{partition}\\ne satisfait les trois};

                    % ---------- Flèche de passage ----------
                    \draw[-{Latex[length=3.2mm]}, line width=3pt, color=gray!40] (4.90, 2.15) -- (6.40, 2.15);
                    \node[font=\scriptsize, align=center, text=gris_fonce_inria] at (5.65, 2.68) {couper\\l'arbre};

                    % ---------- Panneau droit : le dendrogramme du Single-Linkage ----------
                    \begin{scope}[xshift=6.95cm, yshift=0cm]
                        \tikzset{branche/.style={inria-rouge, line width=1.4pt, rounded corners=3pt}}
                        \foreach \x in {0.35, 0.95, 1.55, 2.35, 2.95, 3.95, 4.55, 5.15} {
                            \fill[inria-2024-bleu-canard] (\x, 0.30) circle (2.0pt);
                        }
                        \draw[branche] (0.35,0.30) -- (0.35,0.95) -- (0.95,0.95) -- (0.95,0.30);
                        \draw[branche] (0.65,0.95) -- (0.65,1.45) -- (1.55,1.45) -- (1.55,0.30);
                        \draw[branche] (2.35,0.30) -- (2.35,1.15) -- (2.95,1.15) -- (2.95,0.30);
                        \draw[branche] (3.95,0.30) -- (3.95,0.80) -- (4.55,0.80) -- (4.55,0.30);
                        \draw[branche] (4.25,0.80) -- (4.25,1.60) -- (5.15,1.60) -- (5.15,0.30);
                        \draw[branche] (1.10,1.45) -- (1.10,2.30) -- (2.65,2.30) -- (2.65,1.15);
                        \draw[branche] (1.87,2.30) -- (1.87,3.35) -- (4.70,3.35) -- (4.70,1.60);
                        \draw[branche] (3.28,3.35) -- (3.28,3.95);

                        \draw[coupe] (-0.25, 0.55) -- (5.75, 0.55);
                        \draw[coupe] (-0.25, 1.88) -- (5.75, 1.88);
                        \draw[coupe] (-0.25, 2.85) -- (5.75, 2.85);

                        \draw[decorate, decoration={brace, amplitude=4pt}, inria-rouge, line width=0.8pt]
                            (5.95, 3.05) -- (5.95, 0.35);
                        \node[etiq, anchor=west, align=left] at (6.15, 1.70)
                            {\textsf{(i)}\,+\,\textsf{(ii)}\\ \textsf{(i)}\,+\,\textsf{(iii)}\\ \textsf{(ii)}\,+\,\textsf{(iii)}};

                        \node[font=\scriptsize, text=gris_fonce_inria, anchor=north] at (2.9, -0.15)
                            {selon la \textbf{règle de coupe}, on retrouve chacune des trois paires d'axiomes};
                    \end{scope}
                \end{tikzpicture}
